\section*{Спільні припущення та підготовка}

\markright{Підготовка}

Опціональні пункти 2.б, 2.з, 2.и пропущено.

Дані Ames Housing: 2930 продажів, SalePrice у доларах США, площі у
квадратних футах. Випадкове розбиття 80/20 із seed 42 дає 2344
навчальних і 586 відкладених рядків, однакових для всіх частин; у
частині 2 відкладені рядки називаємо валідаційними, бо на них порівнюємо
кілька моделей. Рівень значущості 0.05.

У частинах 1 і 2 вилучаємо Order, PID і GarageYrBlt; MSSubClass і
текстові ознаки кодуємо індикаторами з одним базовим рівнем, структурне
NA є окремою категорією, числові пропуски заповнюємо навчальною
медіаною, числові ознаки стандартизуємо. Усі перетворення навчаються
лише на навчальній частині, штрафи вибираємо 5-блочною крос-валідацією.
У частині 3 клас 1 означає SalePrice понад 130000 доларів;
використовуємо лише числові ознаки, їх підготовку показано на початку
частини.


\pythoncode{common-data}

\Needspace{120pt}\par\noindent\begin{minipage}[t]{0.49\linewidth}
\begingroup\fontsize{7.5}{9.5}\selectfont\setlength{\tabcolsep}{2.5mm}\renewcommand{\arraystretch}{1.12}\par\noindent\begin{tabular}{|ll|}
\hline
 & Значення \\
\hline
Спостережень & 2930 \\
Навчальна частина & 2344 \\
Перевірочна частина & 586 \\
Числових ознак & 34 \\
Категоріальних ознак & 44 \\
Пропусків у Sale\allowbreak{}Price & 0 \\
\hline
\end{tabular}
\par\endgroup
\end{minipage}\hfill\begin{minipage}[t]{0.49\linewidth}
\begingroup\fontsize{7.5}{9.5}\selectfont\setlength{\tabcolsep}{2.5mm}\renewcommand{\arraystretch}{1.12}\par\noindent\begin{tabular}{|ll|}
\hline
Fireplace\allowbreak{}Qu & Кількість \\
\hline
NA & 1422 \\
Gd & 744 \\
TA & 600 \\
Fa & 75 \\
Po & 46 \\
Ex & 43 \\
\hline
\end{tabular}
\par\endgroup
\end{minipage}\par